\documentclass{resume}

\newcommand{\en}[1]{#1}
\newcommand{\zh}[1]{}

\zh{\usepackage{xeCJK}}
% \zh{\setCJKmainfont{SourceHanSansCN-Regular}}
% \zh{\setCJKsansfont{SourceHanSansHWSC-Regular}}
% \zh{\setCJKmonofont{SourceHanSansHWSC-Regular}}

\begin{document}

\name{\en{Yiwei Yang}\zh{杨易为}}
\basicInfo{
      \email{victoryang00@ucsc.edu} \textperiodcentered\
      \homepage[asplos.dev]{https://asplos.dev}
      \github[victoryang00]{https://github.com/victoryang00}
}


\section{\en{Education}\zh{教育经历}}
\en{\datedsubsection{\textbf{UC Santa Cruz}, Ph.D. Student}{08/2022 -- 06/2028}}
\zh{\datedsubsection{\textbf{加州大学圣克鲁兹分校}, 在读Ph.D.}{2022/09 -- 2028/06}}
\begin{itemize}
      \item \en{Major: Computer Science, advised by \href{https://arquinn.github.io/}{Andrew Quinn}. 22 Fall TA of Computer Architecture}
            \zh{计算机科学，导师为\href{https://arquinn.github.io/}{Andrew Quinn}，2022秋体系结构助教}
\end{itemize}
\en{\datedsubsection{\textbf{ShanghaiTech University}, Undergraduate}{09/2018 -- 06/2022}}
\zh{\datedsubsection{\textbf{上海科技大学}, 本科}{2018/09 -- 2022/06}}
\begin{itemize}
      \item \en{Major: Computer Science, finished Compiler, Network, Database, OS, CA, Convex, RL, Parallel Computing. 21\& 22 Spring TA of Compiler}
            \zh{计算机科学，完成编译原理，操作系统，网络，数据库，并行计算，凸优化，强化学习等课程，2021\&22春编译原理助教}
\end{itemize}


\section{\en{Work Experience}\zh{工作经历}}
\en{\datedsubsection{\textbf{\href{https://www.jumptrading.com/}{Jump Trading}}, Shanghai, China}{07/2020 -- 09/2020}}
\zh{\datedsubsection{\textbf{\href{https://www.jumptrading.com/}{跃申有限公司}}}{2020/07 -- 2020/09}}
\en{\role{Linux Team}{Production Engineer Intern}}
\zh{\role{Linux 组}{HPC自动化运维实习}}
\begin{itemize}
      \item \en{High Frequency Trade Order Book simulation applying Linear.Regression Method.}
            \zh{应用线性回归法进行高频交易盘账模拟。}
      \item \en{Designed a user interface to automate the core affinity of jobs.}
            \zh{在Linux中自动调度作业集成分配cpu核心绑定功能。}
      \item \en{Designed an eBPF exporter of GPFS full OSS lifetime traces for better reporting the bottleneck.}
            \zh{应用eBPF的gobinding对GPFS进行io侦测以更好报系统瓶颈。}
\end{itemize}
\section{\en{Research Experience}\zh{研究经历}}

\en{\datedsubsection{\textbf{\href{https://www.ssrc.ucsc.edu/index.html}{Storage Systems Research Center}, UC Santa Cruz}}{08/2022 -- Present}}
\zh{\datedsubsection{\textbf{\href{https://www.ssrc.ucsc.edu/index.html}{加州大学圣克鲁兹分校存储中心}}}{08/2022 -- 现在}}
\en{\role{Graduate Research}{Assistant}}
\zh{\role{研究生科研}{助理}}
\begin{itemize}
      \item \small{ \en{Data-driven CXL.mem allocation, prefetching, and replacement policies. } \zh{理解CXL.mem的性能分析。数据驱动的远端内存分配、预取和替换策略。}}\includegraphics[width=\linewidth/40]{cxl.png}
      \item \small{ \en{Make Hardware Software Co-design for on CXL.cache data movement}
      \zh{设计一种硬件软件协同的CXL.cache数据迁移方法。}}\includegraphics[width=\linewidth/40]{cxlmemuring.png}
      \item \small{ \en{Make Virtual Machine migration based on WebAssembly and WASI}
      \zh{设计一种基于WebAssembly和WASI的虚拟机迁移方法。}}\includegraphics[width=\linewidth/40]{webassembly.png}
      \item \small{ \en{Make cross kernel-userspace eBPF observation for distributed system}
      \zh{设计一种基于用户态内核态eBPF协同工作的分布式系统。}}\includegraphics[width=\linewidth/40]{ebpf.png}
\end{itemize}

\section{\en{Publications}\zh{论文}}
\href{https://arxiv.org/pdf/2410.15894}{\textbf{"Transparent and Efficient Live Migration across Heterogeneous Hosts with Wharf"} Preprint} \textbf{Yiwei Yang}, Aibo Hu, Yusheng Zheng, Brian Zhao, Xinqi Zhang, Andrew Quinn

\href{https://arxiv.org/pdf/2408.04856}{\textbf{"wasm-bpf: Streamlining eBPF Deployment in Cloud Environments with WebAssembly"} Preprint} Yusheng Zheng, Tong Yu, \textbf{Yiwei Yang}, Andrew Quinn.

\href{https://arxiv.org/pdf/2312.05531}{\textbf{"KEN: Kernel Extensions using Natural Language"} eBPF24} Yusheng Zheng, \textbf{Yiwei Yang}, Maolin Chen, Andrew Quinn.

\href{https://arxiv.org/pdf/2311.07923}{\textbf{"bpftime: userspace eBPF Runtime for Uprobe, Syscall and Kernel-User Interactions"} Preprint} Yusheng Zheng, Tong Yu, \textbf{Yiwei Yang}, Yanpeng Hu, Xiaozheng Lai, Andrew Quinn.

\href{https://arxiv.org/pdf/2303.06153}{\textbf{"CXLMemSim: A pure software simulated CXL.mem for performance characterization."} Yarch23} \textbf{Yiwei Yang} Pooneh Safayenikoo, Jiacheng Ma, Tanvir Ahmad Khan, Andrew Quinn.

\href{https://arxiv.org/pdf/2103.07633}{\textbf{"Attack as Defense: Characterizing Adversarial Examples using Robustness."} ISSTA21} Zhao, Zhe, Guangke Chen, Jingyi Wang, \textbf{Yiwei Yang}, Fu Song, and Jun Sun. 
      
% \en{\datedsubsection{\textbf{\href{http://mir.cs.illinois.edu/marinov/}{Darko Marinov's Lab}, \href{https://illinois.edu}{University of Illinois at Urbana-Champaign}}}{07/2021 -- 10/2021}}
% \zh{\datedsubsection{\textbf{\href{https://illinois.edu}{伊利诺伊大学厄巴纳-香槟分校}\href{http://mir.cs.illinois.edu/marinov/}{Darko Marinov组}}}{2021/07 -- 2021/10}}
% \en{\role{Research Experience for Undergraduate}{Remote Program}}
% \zh{\role{本科生科研}{暑研项目}}
% \begin{itemize}
%       \item \en{Investigating \href{https://victoryang00.cn/wordpress/2022/08/09/final-report-of-program-analysis/}{Order dependent Flaky Tests} by Dynamic Taint Analysis and fix the underlying concurrent bugs.}
%             \zh{使用动态污点分析在Flaky数据集上进行\href{https://victoryang00.cn/wordpress/2022/08/09/final-report-of-program-analysis/}{数据依赖变量检测}，并修复java写的并行系统中的并发bug。}
% \end{itemize}

% \en{\datedsubsection{\textbf{\href{http://s3l.shanghaitech.edu.cn/}{System} \href{https://toast-lab.gitee.io/}{Lab}, ShanghaiTech University}}{07/2019 -- 06/2020}}
% \zh{\datedsubsection{\textbf{\href{http://s3l.shanghaitech.edu.cn/}{上海科技大学系统}\href{https://toast-lab.gitee.io/}{实验室}}}{2019/07 -- 2021/06}}
% \en{\role{Undergraduate Research}{Intern}}
% \zh{\role{本科生科研}{实习}}
% \begin{itemize}
%       \item \small{ \en{A2D uses the cost of attacking an input for robustness evaluation and identifies those less robust examples as adversarial since less robust examples are easier to attack.}
%       \zh{A2D使用攻击输入的成本来进行鲁棒性评估，并将那些鲁棒性较差的label diff rate这一指标guide生成对抗样本。}}\\
%       Zhao, Zhe, Guangke Chen, Jingyi Wang, \textbf{Yiwei Yang}, Fu Song, and Jun Sun. \textbf{"Attack as Defense: Characterizing Adversarial Examples using Robustness."} ISSTA21  .
%       \item \en{Researching \href{https://github.com/LEAFERx/movable}{MOVE language} in Diem currency source code to protect against Arithmetic Overflow, Timestamp Difference.}
%             \zh{研究Diem源码中的\href{https://github.com/LEAFERx/movable}{MOVE language}，提高抵御时间戳攻击、整数溢出等的安全性。}
%       \item \en{Understanding the \href{https://github.com/victoryang00/pmemable}{real access mechanism} of Memory Mode Optane Memory and XPBuffer by reverse-engineering methods.}
%             \zh{通过逆向工程的相关手段去理解在memory mode 下傲腾内存以及XPBuffer的\href{https://github.com/victoryang00/pmemable}{内存置换算法逻辑}。}
% \end{itemize}


\section{\en{Selected Portfolios}\zh{项目}}
\datedsubsection{\href{https://github.com/eunomia-bpf/bpftime/}{\textbf{bpftime\tiny{https://github.com/eunomia-bpf/bpftime/}}}}{\tiny{\en{Make cross boundary observation for the kernel}\zh{提供跨内核态用户态观测机制的工具}}}

\datedsubsection{\href{https://github.com/Multi-V-VM/MVVM}{\textbf{MVVM\tiny{https://github.com/Multi-V-VM/MVVM}}}}{\tiny{\en{Live migration over WebAssembly with WASI support}\zh{一种基于webassembly指令集和wasi系统调用的live migration}}}

\datedsubsection{\href{https://github.com/SlugLab/Bede-linux}{\textbf{Bede-linux\tiny{https://github.com/SlugLab/Bede-linux}}}}{\tiny{\en{Per Process RSS Node Limit Linux kernel for CXL}\zh{为了CXL设计的Per Process RSS Node Limit}}}

% \datedsubsection{\href{https://github.com/victoryang00/rustc-codegen-everything-gpu}{\textbf{rustc-codegen-everything-gpu\tiny{https://github.com/victoryang00/rustc-codegen-everything-gpu}}}}{\tiny{\en{Rust GPU backend for everything}\zh{Rust spirv， webgpu，amdrocm，nvptx后端，为了实现GPU live migration}}}

\section{\en{Skills}\zh{技能}}
\begin{itemize}[parsep=0.25ex]
      \item \en{\textbf{Programming Languages}:
                  not limited to any specific language,
                  and experienced in Python/C++/Rust,
                  comfortable with Golang/C/Java/Scala/TypeScript (in random order).}
            \zh{\textbf{编程语言}:
                  不局限于特定编程语言，且尤其熟悉 Python/C++/Rust 等，
                  了解 Golang/C/Java/Scala/TypeScript 等。}

      \item \en{\textbf{System}:
                  Specialist in Compiler \& Performance Analysis, familiar with \texttt{LLVM}, \texttt{MLIR}, \texttt{Gem5}, \texttt{WASM}, \texttt{gdb}, \texttt{eBPF}, \texttt{qemu}, \texttt{chisel}, \texttt{Linux mm \& observability subsystem}.}
            \zh{\textbf{系统}:
                  熟悉各种CPU，编译器及操作系统内核的概念与设计，熟悉各种内核性能调优工具，例如\texttt{LLVM}, \texttt{MLIR}, \texttt{Gem5}, \texttt{WASM}, \texttt{gdb}, \texttt{eBPF}, \texttt{qemu}, \texttt{chisel}, \texttt{Linux mm \& observability 子系统}。}

      \item \en{\textbf{Machine Learning}:
                  familiar with general knowledge of machine \& reinforce learning, interested of System for LLM for Sys.}
            \zh{\textbf{机器学习}:
                  熟悉经典机器与强化学习算法，对System for LLM for System有兴趣。}
\end{itemize}

\section{\en{Miscellaneous}\zh{其他}}
\begin{itemize}
      \item \en{Interests: Computer Architecture, Storage System, Formal Methods, etc.}
            \zh{兴趣：高性能计算、体系结构、存储、形式化验证安全等。}
      \item \en{Lead \href{https://hpc.geekpie.club}{GeekPie\_HPC} \textit{Ranked 2}, SC-SCC21. \textit{Ranked 4}, ISC22. Advise Not-Slow-Slug \textit{Ranked 2}, ISC23.} \zh{带领\href{https://hpc.geekpie.club}{GeekPie\_HPC} SCC21 团队第2名，ISC22 团队第4名，指导 Not-Slow-Slug 团队 ISC23 团队第2名。}
      % \item \en{As a member of \href{https://github.com/0x238e/Vchain}{0x238e} \textit{Best award}, Bitrun, Hang Zhou, 2019.}
      %       \zh{带领\href{https://github.com/0x238e/Vchain}{0x238e} Bitrun Hackathon 比赛最佳奖}
      \item \en{Google Summer of Code Mentor for BPFTime}
            \zh{谷歌编程之夏 BPFTime 指导}
\end{itemize}

\end{document}
